\thispagestyle{plain}
\phantomsection
\addcontentsline{toc}{chapter}{中文摘要}
\centerline{\zihao{-3}\heiti 中文摘要}

\linespread{1.4}\zihao{-4}\bigskip

本文围绕 2021 年数学建模 A 题“FAST 主动反射面的形状调节”开展建模研究。针对基准球面到工作抛物面的调节任务，建立“理想抛物面确定 -- 节点位移优化 -- 接收比评估”的一体化模型框架。模型在考虑促动器伸缩范围、主索相邻节点距离变化约束和反射几何规律的前提下，以调节后反射面尽量贴近理想抛物面为目标。

在求解上，将 300 m 口径区域的主索节点离散化，构造加权最小二乘贴合目标函数，并结合约束优化算法求解各节点位移和促动器径向伸缩量。随后基于几何光学反射定律进行离散光线追踪，计算馈源舱有效区域接收比，与基准球面状态进行对比分析。

研究表明：该问题可归结为大规模几何约束优化问题，模型结果能够同时输出题目要求的顶点坐标、节点坐标及促动器量，并以接收比指标验证调节方案有效性。该框架具有良好的扩展性，可进一步引入风载、温度和制造误差，服务于 FAST 主动反射面工程化调控分析。

\bigskip
\noindent{\zihao{4}\heiti 关键词：}
FAST；主动反射面；抛物面拟合；约束优化；接收比
